\documentclass[11pt]{article}

\usepackage[margin=1.05in]{geometry}
\usepackage{amsmath,amssymb,amsthm}
\usepackage{graphicx}
\usepackage[colorlinks=true,linkcolor=blue,citecolor=blue,urlcolor=blue]{hyperref}

\newtheorem{theorem}{Theorem}
\newtheorem{corollary}{Corollary}
\newtheorem{lemma}{Lemma}
\newtheorem{remark}{Remark}

\newcommand{\eps}{\varepsilon}
\newcommand{\bp}{\mathrm{BP}}
\newcommand{\ball}{\mathrm{B}}
\newcommand{\Z}{\mathcal Z}
\newcommand{\Czero}{C_0(K)}
\newcommand{\R}{\mathbb R}

\begin{document}

\begin{center}
{\Large A sharp center consequence for almost band preservers on \(C_0(K)\)}\\[0.5em]
{\normalsize Partial result packet for arXiv:1610.02557}\\[0.25em]
{\normalsize Run: \texttt{fa\_banach\_001}}
\end{center}

\section*{Source and target question}

The source paper is Timur Oikhberg and Pedro Tradacete,
\emph{Almost band preservers}, arXiv:1610.02557 \cite{OT}.
It defines a linear map \(T:X\to X\) on a Banach lattice to be
\(\eps\)-band preserving if
\[
 \sup\{\|\,|Tx|\wedge y\,\|: y\ge 0,\ y\perp x\}\le \eps\|x\|
 \qquad (x\in X).
\]
It also defines \(T\in \eps-\Z(X)\) when there is \(\lambda>0\) such that
for every \(x\in \ball(X)\) there is \(z\in X\) with \(\|z\|\le \eps\) and
\[
 |Tx|\le \lambda |x|+z.
\]
Equivalently, as stated in the source paper, \(T\in \eps-\Z(X)\) iff there is
\(\lambda\ge0\) such that
\[
 \sup_{x\in \ball(X)}\|(|Tx|-\lambda |x|)_+\|\le \eps.
\]

The relevant source passage is:

\begin{quote}
``In general, if \(T\in \eps-\Z(X)\), then \(T\) is \(\eps\)-BP.
We do not know whether the converse implication holds in general.''
\end{quote}

\begin{center}
\includegraphics[width=0.92\linewidth]{figures/open_problem_crop.png}
\end{center}

The source later proves that on compact \(C(K)\), every bounded \(\eps\)-BP
operator is within \(2\eps\) of a band-preserving multiplication operator, and
therefore belongs to \(2\eps-\Z(C(K))\) \cite[Theorem 6.2 and following
remark]{OT}. The result below records the sharper center consequence suggested
by the measure estimate inside that proof, and extends it to locally compact
\(C_0(K)\) spaces.

\section*{Result}

\begin{theorem}[Sharp \(C_0(K)\) center consequence]
\label{thm:c0}
Let \(K\) be a locally compact Hausdorff space and let \(\eps\ge0\). If
\[
 T:C_0(K)\longrightarrow C_0(K)
\]
is a linear \(\eps\)-band preserving map, then \(T\in \eps-\Z(C_0(K))\).
More precisely, for every \(f\in C_0(K)\) with \(\|f\|_\infty\le1\),
\[
 \big\| (|Tf|-\|T\|\,|f|)_+\big\|_\infty\le \eps .
\]
\end{theorem}

\begin{corollary}[AL-spaces]
\label{cor:al}
Let \(X\) be an AL-space and let \(\eps\ge0\). Every linear
\(\eps\)-band preserving map \(T:X\to X\) belongs to \(\eps-\Z(X)\).
\end{corollary}

\begin{remark}[Sharpness]
\label{rem:sharp}
The constant \(\eps\) in Theorem \ref{thm:c0} is optimal, even for
\(\ell_\infty^2\). For \(\eps>0\), define
\[
 T(a,b)=(\eps b,0),\qquad (a,b)\in \ell_\infty^2.
\]
Then \(T\) is \(\eps\)-BP. However, for \(x=(0,1)\) and every
\(\lambda\ge0\),
\[
 (|Tx|-\lambda |x|)_+=(\eps,0),
\]
so \(T\notin c-\Z(\ell_\infty^2)\) for any \(c<\eps\).
\end{remark}

\section*{Proof intuition}

For \(C_0(K)\), boundedness is automatic by the source paper. At each
point \(t\in K\), the scalar functional \(f\mapsto Tf(t)\) is represented by
a regular Radon measure \(\mu_t\). Split off the atom at \(t\):
\[
 \mu_t=a_t\delta_t+\nu_t,\qquad \nu_t(\{t\})=0.
\]
The \(\eps\)-BP condition says that a function vanishing at \(t\) cannot
produce more than \(\eps\) at \(t\) after applying \(T\). This controls the
off-atom part: \(\|\nu_t\|\le\eps\). Hence
\[
 |Tf(t)|\le |a_t|\,|f(t)|+\eps\|f\|_\infty,
\]
and \(|a_t|\le\|T\|\). That pointwise inequality is exactly the
\(\eps\)-center inequality.

\section*{Proof}

We first record the locally compact version of the scalar vanishing criterion
used in the compact \(C(K)\) proof of \cite[Lemma 6.1]{OT}.

\begin{lemma}
\label{lem:vanish-at-point}
Let \(T:C_0(K)\to C_0(K)\) be linear and \(\eps\)-BP. Then \(T\) is bounded,
and for every \(t\in K\) and every \(f\in C_0(K)\) with \(f(t)=0\),
\[
 |Tf(t)|\le \eps\|f\|_\infty .
\]
\end{lemma}

\begin{proof}
The boundedness is \cite[Theorem 3.13]{OT}, applied to the scalar space
\(\mathbb R\).

If \(f\) vanishes on an open neighborhood \(V\) of \(t\), the desired estimate
is exactly \cite[Lemma 3.9]{OT}, again in the scalar case.

Now suppose only that \(f(t)=0\). Fix \(\delta>0\). By continuity, choose an
open neighborhood \(U\) of \(t\) such that \(|f(s)|<\delta\) on \(U\). Since
\(K\) is locally compact Hausdorff, choose \(h\in C_c(K)\) with
\(0\le h\le1\), \(h=1\) on a neighborhood of \(t\), and
\(\operatorname{supp}h\subset U\). Put \(g=(1-h)f\). Then \(g\) vanishes on
a neighborhood of \(t\), \(\|g\|_\infty\le\|f\|_\infty\), and
\(\|f-g\|_\infty=\|hf\|_\infty\le\delta\). Therefore
\[
 |Tg(t)|\le \eps\|f\|_\infty,
 \qquad
 |T(f-g)(t)|\le \|\delta_t T\|\,\delta .
\]
Letting \(\delta\downarrow0\) gives the assertion.
\end{proof}

\begin{proof}[Proof of Theorem \ref{thm:c0}]
By Lemma \ref{lem:vanish-at-point}, \(T\) is bounded. Fix \(t\in K\). The
bounded functional \(f\mapsto Tf(t)\) on \(C_0(K)\) is represented by a
regular finite signed (or complex) Radon measure \(\mu_t\):
\[
 Tf(t)=\int_K f\,d\mu_t,\qquad f\in C_0(K).
\]
Let \(a_t=\mu_t(\{t\})\) and \(\nu_t=\mu_t-a_t\delta_t\). Then
\(\nu_t(\{t\})=0\).

We claim that \(\|\nu_t\|\le\eps\). Let \(f\in C_0(K)\) with
\(\|f\|_\infty\le1\), and let \(\eta>0\). By regularity of the total variation
measure \(|\nu_t|\), and since \(|\nu_t|(\{t\})=0\), choose an open
neighborhood \(U\) of \(t\) such that \(|\nu_t|(U)<\eta\). Choose
\(h\in C_c(K)\) with \(0\le h\le1\), \(h(t)=1\), and
\(\operatorname{supp}h\subset U\). Then
\[
 |\nu_t(fh)|\le |\nu_t|(U)<\eta.
\]
Also \(f(1-h)\) vanishes at \(t\), so Lemma \ref{lem:vanish-at-point} gives
\[
 |\mu_t(f(1-h))|=|T(f(1-h))(t)|\le\eps .
\]
Because \(\delta_t(f(1-h))=0\), this is the same as
\(|\nu_t(f(1-h))|\le\eps\). Hence
\[
 |\nu_t(f)|\le \eps+\eta.
\]
Letting \(\eta\downarrow0\) and taking the supremum over the unit ball of
\(C_0(K)\) proves \(\|\nu_t\|\le\eps\).

For \(\|f\|_\infty\le1\), we therefore have pointwise
\[
 |Tf(t)|=|a_t f(t)+\nu_t(f)|
 \le |a_t|\,|f(t)|+\eps .
\]
Moreover \(|a_t|\le\|\mu_t\|=\|\delta_t T\|\le\|T\|\). Thus
\[
 |Tf(t)|\le \|T\|\,|f(t)|+\eps
 \qquad (t\in K,\ \|f\|_\infty\le1),
\]
which is equivalent to
\[
 \big\|(|Tf|-\|T\|\,|f|)_+\big\|_\infty\le\eps.
\]
This is precisely \(T\in\eps-\Z(C_0(K))\).
\end{proof}

\begin{proof}[Proof of Corollary \ref{cor:al}]
For \(\eps>0\), an \(\eps\)-BP linear map on an order continuous Banach
lattice is bounded by \cite[Corollary 3.8]{OT}; AL-spaces are order
continuous. For \(\eps=0\), boundedness of band-preserving maps is the
Abramovich--Veksler--Koldunov theorem quoted in the source introduction.
Thus \(T\in B(X)\).

Since \(X\) is order continuous, \cite[Proposition 2.11]{OT} implies that
\(T^*:X^*\to X^*\) is \(\eps\)-BP. The dual of an AL-space is an AM-space
with unit, hence is lattice isometric to \(C(L)\) for a compact Hausdorff
space \(L\) by Kakutani's representation theorem. Applying Theorem
\ref{thm:c0} to this compact \(C(L)\) gives
\[
 T^*\in \eps-\Z(X^*).
\]
Finally, \cite[Proposition 4.3]{OT} says that \(S\in\eps-\Z(Y)\) iff
\(S^*\in\eps-\Z(Y^*)\). Applying the reverse direction to \(T\) gives
\(T\in\eps-\Z(X)\).
\end{proof}

\section*{Verification notes and limitations}

The proof is a locally compact adaptation of the measure estimate already
present in \cite[Theorem 6.2]{OT}. The source paper used this estimate to
approximate \(T\) by the multiplication operator \(f\mapsto (T1)f\), losing a
second \(\eps\). For the center conclusion that approximation step is
unnecessary: the off-atom estimate alone gives the sharp \(\eps\)-bound.

This packet does not solve the source paper's general question for arbitrary
Banach lattices. It gives the sharp positive answer for \(C_0(K)\) spaces and,
by duality, for AL-spaces. Local run-index searches for the arXiv id and core
phrases found no duplicate packet. A bounded web search for exact title and
keyword combinations found the source arXiv record and no later exact
resolution of this sharpened statement.

\begin{thebibliography}{9}

\bibitem{OT}
T. Oikhberg and P. Tradacete,
\emph{Almost band preservers},
arXiv:1610.02557, 2016.
\url{https://arxiv.org/abs/1610.02557}

\end{thebibliography}

\end{document}
